\documentclass[11pt,a4paper]{article}
\usepackage[margin=2cm]{geometry}
\usepackage[T1]{fontenc}
\usepackage{mathptmx}
\usepackage{amsmath,amssymb}
\usepackage{xcolor}
\usepackage{enumitem}
\usepackage{array}
\usepackage[hidelinks]{hyperref}
\definecolor{navy}{HTML}{1B3755}
\definecolor{blue}{HTML}{2F6FB3}
\setlength{\parindent}{0pt}
\setlength{\parskip}{5pt}
\renewcommand{\arraystretch}{1.25}
\newcommand{\sect}[1]{\bigskip{\Large\color{navy}\bfseries #1}\par\smallskip{\color{navy}\hrule height 0.8pt}\smallskip}
\newcommand{\say}[1]{\par\smallskip\noindent\fbox{\parbox{\dimexpr\linewidth-2\fboxsep-2\fboxrule}{\textbf{How to say it:} #1}}\par}

\begin{document}
\begin{center}
{\LARGE\bfseries\color{navy} Fuse Revision in the DSDR Study}\\[4pt]
{\large What was changed, when, why and how}\\[2pt]
IEEE 13-node feeder, DIgSILENT PowerFactory --- Jhala Nath Kafle (081MSPSE009)
\end{center}

\sect{1. What --- the fuse sizes before and after}
All fuses are A055C (E-rated) library fuses, except F634 (gL-800A, 0.48\,kV side). Only the fuses below were
changed; the other eight kept their starting size.

\begin{center}
\begin{tabular}{|l|l|c|c|c|}
\hline
\textbf{Fuse} & \textbf{Location} & \textbf{Start} & \textbf{Stage A (no DG)} & \textbf{Stage C (DG + DSDR)} \\ \hline
F692-R & cable 692--675, at 692 & 250E & \textbf{200E} & \textbf{250E} \\ \hline
F-DL & distributed-load lateral on 632--671 & 250E & 250E & \textbf{400E} \\ \hline
F633 & lateral 632--633, at 632 & 250E & 250E & \textbf{400E} \\ \hline
F671-1 & 671--692 (switch), at 671 & 300E & 300E & \textbf{400E} \\ \hline
F646 & line 645--646, at 645 & 300E & 300E & \textbf{400E} \\ \hline
F632 & lateral 632--645, at 632 & 400E & 400E & \textbf{500E} \\ \hline
\end{tabular}
\end{center}

\sect{2. When --- two moments in the method}
The fuse size is the \emph{last} thing the method changes. The order is always: settings first, then the
time dials (step 8), and only when a dial is at its limit, the fuse size (step 9) and the series check (step 10).
\begin{enumerate}[leftmargin=*,itemsep=3pt]
\item \textbf{Stage A, design without DG (step 10, series fuses).} One change: F692-R 250E $\rightarrow$ 200E.
\item \textbf{Stage C, DG connected and R2 as a DSDR (step 9 and 10).} After the dial revision (step 8) could not
  improve any lost case. Six fuses changed.
\end{enumerate}
Stage B (DG connected, conventional R2) only \emph{evaluates} the old design; nothing is changed there.

\sect{3. Why --- what was failing}
\textbf{Stage A (no DG).} With the starting sizes 35 of 39 cells were held. The four lost cells were all at 675,
and the reason was the series-fuse condition between F692-R (250E) and its upstream fuse F671-1 (300E): F692-R did
not clear early enough. Example, LLL at 675: F692-R clears in 0.293\,s, but the limit set by F671-1's melting
curve is 0.198\,s. A smaller F692-R (200E) clears faster $\rightarrow$ 39 of 39.

\textbf{Stage C (DG + DSDR).} With the DG the fuse carries the grid current \emph{plus} the DG current, while the
recloser carries only the grid share. The fuse therefore melts before the recloser's fast trip. Examples from the
evaluation with DG (old fuses):

\begin{center}
\begin{tabular}{|l|l|c|c|}
\hline
\textbf{Fault} & \textbf{Fuse} & \textbf{Fuse starts to melt} & \textbf{Recloser fast trip} \\ \hline
633, LLL & F633 (250E) & 0.055\,s & R1: 0.106\,s \\ \hline
DL, LLL & F-DL (250E) & 0.046\,s & R1: 0.106\,s \\ \hline
646, LL (b--c) & F646 (300E) & 0.178\,s & R1: 0.189\,s \\ \hline
675, LLL & F692-R (200E) & 0.038\,s & R2: 0.059\,s \\ \hline
\end{tabular}
\end{center}

\textbf{Why the dials could not fix it (step 8).} To win against the fuse the recloser's fast curve would have to
be faster. But R1's fast dial is already at the bottom of the IAC range (TDS 0.5) and R2's at the bottom of the CDG
range (TMS 0.1). The delayed dials are at their tops (TDS 10, TMS 1.0). So the only remaining step of the method
is a \emph{slower fuse}, i.e.\ a larger size.

\textbf{Why some backup fuses had to grow as well.} When a fuse grows, the fuse upstream of it (the backup) must
still be slower, otherwise the series-fuse condition (eq.~8) fails. So F692-R up needs F671-1 up, and F646 up needs
F632 up.

\sect{4. How --- the procedure}
\begin{enumerate}[leftmargin=*,itemsep=3pt]
\item Classify all 39 cells (every fault type and phase combination) with the current settings, with the DG
  connected \emph{and} disconnected, and list the lost fault cases with their reason.
\item From each lost case build candidate changes:
  \begin{itemize}[itemsep=1pt]
  \item fuse melts before a fast trip $\rightarrow$ that fuse 1, 2 or 3 A055C sizes \emph{up}, alone or together
    with its backup fuses;
  \item fuse clears after a delayed trip $\rightarrow$ that fuse 1--3 sizes \emph{down};
  \item series condition fails $\rightarrow$ backup fuse up, or primary fuse down.
  \end{itemize}
\item Try every candidate on the whole feeder (all cells, DG in and out) and keep the one that removes the most
  lost cases with the smallest change.
\item Apply it and repeat from step 1 until no candidate improves the result.
\end{enumerate}
Only real A055C sizes were used, so every result is a fuse that can be bought.

\textbf{The changes in the order they were applied (stage C).} The count is of lost fault \emph{cases} (each phase
combination counts once, with the DG in and out), not of cells.

\begin{center}
\begin{tabular}{|c|>{\raggedright}p{4.3cm}|c|>{\raggedright\arraybackslash}p{7.0cm}|}
\hline
\textbf{Order} & \textbf{Change} & \textbf{Lost cases} & \textbf{Reason} \\ \hline
1 & F-DL 250E $\rightarrow$ 400E & 30 $\rightarrow$ 20 & F-DL melted before R1's fast trip for every fault at DL \\ \hline
2 & F633 250E $\rightarrow$ 400E & 20 $\rightarrow$ 10 & F633 melted before R1's fast trip for faults at 633 \\ \hline
3 & F692-R 200E $\rightarrow$ 250E \emph{and} F671-1 300E $\rightarrow$ 400E & 10 $\rightarrow$ 9 &
  F692-R melted before R2's fast trip at 675; its backup F671-1 grows with it (series condition) \\ \hline
4 & F646 300E $\rightarrow$ 400E \emph{and} F632 400E $\rightarrow$ 500E & 9 $\rightarrow$ 8 &
  F646 melted before R1's fast trip at 646; its backup F632 grows with it \\ \hline
\end{tabular}
\end{center}
F692-R returns to 250E in stage C: with the DG it must be slower than R2's fast trip, and the larger F671-1
(400E) now leaves room for it in the series condition.

\sect{5. Result --- how much the revision gives}
\begin{center}
\begin{tabular}{|l|c|c|}
\hline
\textbf{DG connected} & \textbf{Fuses of the no-DG design} & \textbf{Revised fuses} \\ \hline
R2 conventional (single setting) & 24 & 31 \\ \hline
R2 as DSDR (dual setting) & 25 & \textbf{38} \\ \hline
\end{tabular}
\end{center}
\begin{itemize}[leftmargin=*,itemsep=2pt]
\item The DSDR alone: 24 $\rightarrow$ 25. The fuse revision alone: 24 $\rightarrow$ 31. Both: 24 $\rightarrow$ 38.
\item So the two measures work together: seven of the fourteen restored cells need the dual setting
  (633 all types; LG at 645, 646 and DL), the others need the larger fuses.
\item With the final settings and the DG \emph{disconnected}, 37 of 39 cells are held.
\end{itemize}

\textbf{What the revision could not fix --- LG fault at 692.} F671-1 must be 400E to stay selective with F692-R.
For the small LG current at 692 the 400E fuse is slow: it clears in about 7.6\,s, while R2's delayed curve trips at
about 4.0\,s. R2's delayed dial is already at its maximum (TMS 1.0, the highest that keeps R2 below R1), so the
cell stays lost.

\sect{6. Side effects that were checked}
\begin{itemize}[leftmargin=*,itemsep=3pt]
\item \textbf{Slower clearing of weak faults.} A larger fuse is slower at small currents. That is why two of the
  nine selected faults (cases 3 and 4, faults through impedance) are lost with both settings.
\item \textbf{Sympathetic operation of F671-1.} With the DG, F671-1 carries DG current for faults outside its zone.
  For every such fault R2's reverse fast trip comes first, e.g.\ fault at 671 (LLL): F671-1 would melt at 2.467\,s,
  R2 trips at 0.054\,s.
\item \textbf{Load.} Every revised fuse is larger than before, so it still carries its load current.
\end{itemize}

\sect{7. Short answers for the defence}
\say{``The fuse revision is step 9 of the method. It is used only after the dial revision fails, and here the dials
were already at their limits: R1 at TDS 0.5 and R2 at TMS 0.1. With the DG the fuse carries grid plus DG current
and melts before the fast trip, so the fuse had to be made slower, which means one or two sizes larger.''}
\say{``Six fuses changed: F-DL, F633, F646 and F671-1 to 400E, F632 to 500E and F692-R back to 250E. Each change
was chosen by trying every candidate on the whole feeder and keeping the one that removed the most lost cases.''}
\say{``Is it cheating? No: the revision alone gives 31 cells; the dual setting is still needed for seven more,
giving 38.''}
\end{document}
